\documentclass[11pt]{article}

\usepackage[utf8]{inputenc}
\usepackage[T1]{fontenc}
\usepackage{lmodern}
\usepackage{geometry}
\usepackage{amsmath,amssymb,amsfonts,amsthm,mathtools}
\usepackage{enumitem}
\usepackage{microtype}
\usepackage{hyperref}

\geometry{margin=1in}
\hypersetup{colorlinks=true, linkcolor=black, urlcolor=blue, citecolor=black}
\setlist[enumerate]{topsep=0.35em,itemsep=0.3em}
\setlist[itemize]{topsep=0.3em,itemsep=0.25em}

\newcommand{\Aether}{\AE{}ther}
\newcommand{\Aflow}{\AE{}ther-flow}
\newcommand{\TheoryProgram}{\Aether{} / \Aflow{} framework}
\newcommand{\Sub}{\mathcal{A}}
\newcommand{\PhiFlow}{\Phi_{\lambda}}
\newcommand{\Step}{\mathsf{Step}}
\newcommand{\Adj}{\mathsf{Adj}_{\mathcal{S}}}
\newcommand{\RespAlg}{\mathcal{R}_{\Gamma}}
\newcommand{\Rank}{\mathsf{rank}_{\mathcal{S}}}
\newcommand{\Cycle}{\mathcal{C}_{\Gamma}}
\newcommand{\Basis}{\mathcal{B}_{\mathrm{src}}}
\newcommand{\Lang}{\mathcal{L}_{0}}
\newcommand{\Ext}{\mathcal{E}_{\Basis}}

\newtheorem{axiompacket}{Axiom packet}
\newtheorem{justification}{Justification}
\newtheorem{classification}{Conservativity classification}
\newtheorem{auditburden}{Audit burden}
\newtheorem{localverdict}{Local verdict}
\newtheorem{nextstep}{Next step}

\title{\textbf{Localization Source-Basis Axiom Justification}\\[0.35em]
\large Minimal Source-Side Extension Hypotheses After the Conservativity Refutation}
\author{Ontology Formalizer AgentJob AJ-RT-20260608-010-001}
\date{June 8, 2026}

\begin{document}

\maketitle

\begin{abstract}
This draft responds to the local source-basis conservativity refutation. The
result is not a derivation and not a candidate reconstruction. It proposes the
minimum current source-side axiom packet that would make the transition,
adjacency, perturbation, response-rank, and clock-transport structures
determinately testable. The conclusion is restrained: these axioms are not
conservative over the original registered source reservoir. They are explicit
ontology-extension hypotheses that may be audited, refuted, or later promoted
only through a separate authorized gate.
\end{abstract}

\section{Scope and Standard}

This artifact is a draft/control output for task
\texttt{RT-20260608-010}. It does not edit canonical ontology TeX, promote a
localization candidate, issue a Gate Chair verdict, or claim completion of an
effective relativistic sector.

Let \(\Lang\) denote the registered source language available before the
source-basis packet: the substrate arena, intrinsic ordered motion, observer
slice placeholder, S-time placeholder, and the downstream draft localization
control artifacts. The previous Refuter result showed that the additional
source-basis data
\[
\Ext =
(\Step,\Adj,\mathcal{P},\mathcal{D},\{\RespAlg\}_{\Gamma},\mathcal{K})
\]
are not forced by \(\Lang\). The present task therefore cannot honestly
declare them conservative over \(\Lang\). It can only do one of two things:
\begin{enumerate}
    \item identify a conservative definition relative to already stated
    extension data; or
    \item mark a structure as a new source-side axiom whose adoption would be
    an ontology-extension hypothesis.
\end{enumerate}

\section{Minimal Axiom Packet}

\begin{axiompacket}[A0: finite source presentation]
For a localization test, the substrate arena is represented by a finite or
locally finite source presentation
\[
(\Sub_U,\preceq_U,\Phi_U)
\]
on a bounded source-accessible domain \(U\). This presentation is an
experimental/control representation of the substrate order, not a target
manifold chart. Relabelings of \(\Sub_U\) and monotone reparameterizations of
the order remain presentation redundancy.
\end{axiompacket}

\begin{justification}
The original source language is too spare to determine localization-relevant
incidence. A locally finite presentation is the weakest current control
device that lets the project ask whether transition, response, and transport
claims are invariant under relabeling. It is not claimed as canonical
ontology.
\end{justification}

\begin{axiompacket}[A1: source transition germ]
The immediate transition relation \(\Step\) is the primitive edge relation of
the locally finite source presentation. It must satisfy:
\begin{enumerate}
    \item order compatibility: \(x\Step y\) implies \(x\preceq_U y\) or
    \(y\preceq_U x\) according to the directed presentation convention;
    \item locality: every edge has source support inside the bounded domain
    used by the perturbation test;
    \item relabeling invariance: if a source relabeling preserves
    \((\Sub_U,\preceq_U,\Phi_U)\), it preserves the edge-incidence class of
    \(\Step\); and
    \item no benchmark fitting: \(\Step\) is fixed before any comparison with
    target causal, spatial, or clock data.
\end{enumerate}
\end{axiompacket}

\begin{classification}[Transition status]
\(\Adj\), neighborhoods, and source boundaries may be conservative
definitions after A1, because they can be generated from \(\Step\). A1 itself
is not conservative over \(\Lang\). It is an explicit source-transition
extension hypothesis.
\end{classification}

\begin{axiompacket}[A2: admissible perturbations and discriminators]
An admissible perturbation class \(P_a\in\mathcal{P}\) is a finite
source-local change to \(\Step\) and source response records near \(a\) that
preserves all declared source constraints outside its support. A discriminator
\(D\in\mathcal{D}\) is admissible only when it is fixed by source recurrence
words, transition counts, and chain-internal response distinctions before
benchmark comparison.
\end{axiompacket}

\begin{justification}
The Refuter identified \((P_a,D)\) as a discrete tuning channel. A2 blocks
post hoc threshold selection by requiring every detector distinction to be
specified from source recurrence structure. If no such distinction is
available, the candidate must report non-detection or non-uniqueness.
\end{justification}

\begin{classification}[Perturbation status]
The rank \(r_D(a,b)\) is a conservative definition once A1 and A2 are adopted.
A2 is not conservative over \(\Lang\), because the original source reservoir
does not determine perturbation classes or discriminator families.
\end{classification}

\begin{axiompacket}[A3: response algebra without fixed target rank]
For each observer chain \(\Gamma\), \(\RespAlg\) is the quotient algebra of
source response words generated by concatenation, recurrence counting, and
source-local recombination, modulo relabeling symmetries of the source
presentation. The rank \(\Rank(U_{\Gamma})\) is the independence rank of
response functionals in this quotient. No axiom sets \(\Rank=4\).
\end{axiompacket}

\begin{justification}
The previous packet risked selecting the target \(1+3\) split by requiring a
rank-four domain. A3 reverses the burden. The source algebra reports its rank;
a later candidate may use a four-rank domain only if the source records supply
one without post hoc selection.
\end{justification}

\begin{classification}[Response-rank status]
Rank computation is conservative relative to a fixed quotient response
algebra. The quotient rule and admissible response-word ground set are
ontology-extension hypotheses unless they are later derived from canonical
source dynamics.
\end{classification}

\begin{axiompacket}[A4: clock-cycle transport as source holonomy data]
A stable cycle class \([\Cycle]_{\mathcal{K}}\) is a source recurrence class
whose transition-word and response-ratio records remain bounded under
iteration. A transport map
\[
T_{\Gamma\to\Gamma'}:
[\Cycle]_{\mathcal{K}}\to[\Cycle']_{\mathcal{K}}
\]
is admissible only when induced by shared source transition support between
adjacent chains. Transport around closed paths may have nontrivial holonomy;
the holonomy must be reported as source data, not normalized away to match a
target clock.
\end{axiompacket}

\begin{justification}
The previous clock-scale repair added transport but did not determine it.
A4 fixes the anti-smuggling rule: clock comparison is source-holonomy data
first and benchmark clock comparison later. Trivial holonomy is not assumed.
\end{justification}

\begin{classification}[Clock-transport status]
Cycle equivalence and transport are conservative definitions only after the
shared transition-support relation and holonomy convention are fixed. Those
fixing rules are ontology-extension hypotheses over \(\Lang\).
\end{classification}

\section{Conservativity Ledger}

\begin{center}
\begin{tabular}{p{0.25\linewidth}p{0.32\linewidth}p{0.32\linewidth}}
\textbf{Object} & \textbf{Status over original \(\Lang\)} & \textbf{Status after axioms} \\
\hline
\(\Step\) & Nonconservative extension & A1 primitive to audit \\
\(\Adj\), \(N_n(x)\), \(\partial_{\mathcal{S}}R\) & Nonconservative without \(\Step\) & Conservative definitions from A1 \\
\(\mathcal{P}\), \(\mathcal{D}\) & Nonconservative extension & A2 primitives to audit \\
\(r_D(a,b)\) & Underdefined & Conservative definition after A1 and A2 \\
\(\RespAlg\), \(\Rank\) & Nonconservative extension & A3 quotient and rank survey \\
\(\mathcal{K}\), \(T_{\Gamma\to\Gamma'}\) & Nonconservative extension & A4 holonomy data to audit \\
\end{tabular}
\end{center}

The ledger makes the central conclusion explicit. The packet does not rescue
conservativity over the original source reservoir. It records the minimal
current extension hypotheses that would make the next candidate path
determinately falsifiable.

\section{Failure Tests}

\begin{auditburden}[Transition audit]
Reject A1 for localization use if the chosen \(\Step\) is selected to
reproduce target neighborhoods, target causal fronts, or target spatial balls
rather than fixed by source presentation rules.
\end{auditburden}

\begin{auditburden}[Discriminator audit]
Reject A2 if any response threshold, recurrence distinction, or perturbation
class is chosen after benchmark comparison or is adjusted to improve target
agreement.
\end{auditburden}

\begin{auditburden}[Rank audit]
Reject A3 for candidate construction if rank four is imposed as an assumption
instead of reported as a source-rank result. Lower or higher source rank is a
valid negative result.
\end{auditburden}

\begin{auditburden}[Transport audit]
Reject A4 if transport holonomy is normalized to target proper-time behavior
without a source reason. Nontrivial holonomy is not a failure by itself; hiding
it is the failure.
\end{auditburden}

\section{Local Verdict}

\begin{localverdict}[Axioms justified only as extension hypotheses]
The minimal source-basis axiom packet is justified as a controlled
ontology-extension hypothesis, not as a conservative definition over the
existing source language. It improves the research state by converting the
Refuter's nonconservative channels into explicit, auditable assumptions. It
does not authorize candidate reconstruction.
\end{localverdict}

\section{Next Role}

\begin{nextstep}[Smuggling audit before reconstruction]
The logical next role is a Smuggling Auditor review of A1--A4. The audit
should test whether the new source-side axioms secretly encode target
topology, target observer protocols, target rank, or target clock
normalization. Candidate Constructor and Gate Chair review remain premature.
\end{nextstep}

\section*{References}

Ricciardi, A. S. (2026, June 8). \emph{Localization source-basis
conservativity refutation: Source-language test for the source-basis packet}
[Unpublished manuscript]. The Aether-Flow Interpretation of Relativity
Research Project,
\texttt{research\_control/tasks/RT-20260608-009/artifacts/08\_LOCALIZATION\_SOURCE\_BASIS\_CONSERVATIVITY\_REFUTATION.tex}.

Ricciardi, A. S. (2026, June 8). \emph{Localization source-basis smuggling
audit: Target-import review of the source-basis formalization} [Unpublished
manuscript]. The Aether-Flow Interpretation of Relativity Research Project,
\texttt{research\_control/tasks/RT-20260608-008/artifacts/07\_LOCALIZATION\_SOURCE\_BASIS\_SMUGGLING\_AUDIT.tex}.

Ricciardi, A. S. (2026, June 8). \emph{Localization source-basis
formalization: Minimal order, adjacency, response, mode, and clock-transport
structures} [Unpublished manuscript]. The Aether-Flow Interpretation of
Relativity Research Project,
\texttt{research\_control/tasks/RT-20260608-007/artifacts/06\_LOCALIZATION\_SOURCE\_BASIS\_FORMALIZATION.tex}.

Ricciardi, A. S. (2026, June 8). \emph{Localization repair refutation:
Source-determination test for the repaired localization packet} [Unpublished
manuscript]. The Aether-Flow Interpretation of Relativity Research Project,
\texttt{research\_control/tasks/RT-20260608-006/artifacts/05\_LOCALIZATION\_REPAIR\_REFUTATION.tex}.

Ricciardi, A. S. (2026, June 8). \emph{Gate 0 ontology formalization: Source
primitives, forbidden imports, and localization burden} [Unpublished
manuscript]. The Aether-Flow Interpretation of Relativity Research Project,
\texttt{research\_control/tasks/RT-20260608-002/artifacts/01\_GATE0\_ONTOLOGY\_FORMALIZATION.tex}.

\end{document}
